\section{Discussion}\label{sec:discussion}

\subsection{What is established}

Sections~\ref{sec:kernel} and~\ref{sec:frames} are, mathematically, elementary: the complex quaternion form of the Riemann--Silberstein vector,
its Lorentz invariant, its duality orbits and the Stokes picture of the doublet \((E,B)\) are known in substance~\cite{bialynicki}.  What they
add is a precise form for the readout hypothesis: two sectors in one orbit, a null sector fixed by duality, the symmetric triple as the centred
null configuration, and an involution \(J_B=R_{\pi/2}\circ T\) that exchanges the sectors and contains time reversal.  Section~\ref{sec:generation}
is a construction under a stated setting.  Its algebra is largely known --- it is the Pati--Salam structure inside \(\mathbb C\otimes\HH\otimes\OO\)
--- and its specific content is the observation that the projector restricting the weak \(SU(2)\) has compensator form and is singled out by
anomaly cancellation among four candidates.  Nothing in Sections~\ref{sec:kernel}--\ref{sec:frames} enters its proofs; the common element is the
form \(\tfrac12(1+J)\) and the role of time reversal, which Section~\ref{sec:readout} makes explicit.

\subsection{What a dynamics would have to do}

The common gap is the absence of an action varied in all fields.  The readout hypothesis and the results above constrain such an action.
\begin{enumerate}
\item Coframe, connection and matter must be independent fields that act back; a geometry assembled from probe fields alone becomes flat.
\item The positive and the negative reading need the same kinetic term, so that the map between them is a symmetry of the action, and the
hypothesis requires this map to contain time reversal.
\item The parameter \(t\) of the kernel must acquire a kinematic meaning, as a frequency or a mass, so that its imaginary value on the negative
sector (Section~\ref{sec:readout}) has consequences.
\item The compensator must enter as a coupling, not as a projector applied afterwards.
\item \(\Lambda\) should appear as an integration constant, as in unimodular gravity~\cite{henneaux} or vacuum-energy sequestering~\cite{kaloper}.
\item A dimensionless structural value such as \(\eta\) has to come from a dimensionless ratio; a frozen dimensionful value would carry the
expansion history.
\end{enumerate}
A comparison model of probe field equations in which the coframe is reconstructed from the spinor field has been analysed with the formal
theory of partial differential equations~\cite{bryant,seiler}; it will be reported separately.

\subsection{Open questions}

Is the anomaly selection of Proposition~\ref{prop:anomaly} stated in the literature, and does it survive other placements of the weak
\(SU(2)\), for instance the one of~\cite{fureyhughes}?  Is there a realization of \eqref{eq:cm} on a 16-dimensional carrier that satisfies the
spectral hypotheses of Section~\ref{sec:frames}?  Can the step from a time-reversed reading to space-like propagation be made in any
dynamics?  And is there a reason, independent of the lepton data, for the ratio \eqref{eq:gammaL}?

\subsection{Relation to an earlier version}

This paper condenses a much longer, openly developed text~\cite{companion}.  Several statements made there do not hold, among them all steps
that relied on a claimed proof of the Riemann hypothesis, a cosmological-constant argument that is off by a factor of about \(10^{121}\), and
six earlier predictions, none of which is confirmed.  A list with reasons is kept in the repository~\cite{repo}; none of it is used here.
